\texttt{centroids(g)} 返回全部重心，所以结果长度一定是 $1$ 或 $2$；有两个重心时二者必相邻。

\texttt{diameter(g)} 返回 \texttt{\{长度,端点a,端点b\}}，长度按边数计算。
函数连续做两次 BFS：先从任意点找到直径端点 $a$，再从 $a$ 找到另一个端点 $b$。

这两个函数刻意不封装成结构体，比赛时传入 \texttt{vector<vector<int>>} 就能直接取结果。
带权树不能直接使用这里的 BFS 直径函数，应改成 DFS 累加边权；若允许负边权，“任意点两次最远点”结论也不再成立。
